\originalchapter{磁矩、交换与自旋波}
\label{ch:magnet}
\source{8.512第14--19讲，8.513相互作用玻色子组的磁序主题。}
\originalsection{局域磁矩与交换尺度}
原子部分填满的壳层具有自旋与轨道角动量。Hund规则是在一定相互作用与能级顺序下先使总自旋尽可能大，再判断轨道及自旋轨道耦合；晶场可使轨道角动量淬灭，不能在晶体里无条件套孤立原子的简并。独立局域磁矩的小场配分函数展开给$\chi=C/T$，$C=\mu_0n(g\mu_B)^2S(S+1)/(3\kb)$，与Pauli顺磁近常数的温度行为不同。

相邻局域轨道通过虚跃迁获得超交换。本册两站Hubbard推导给$J=4t^2/U>0$。若包含多个轨道及Hund耦合，虚中间态能量不同，交换可改变符号。金属中的局域矩还可通过传导电子产生RKKY作用$J_{ij}\propto J_K^2\chi_0(R_i-R_j)$；三维长距离随$\cos(2k_FR)/R^3$振荡。该结果来自把电子密度--自旋响应积分到二阶，而非单一近邻超交换的延长。
\originalsection{铁磁自旋波}
铁磁Heisenberg模型$H=-J\sum_{\avg{ij}}\vct S_i\cdot\vct S_j$，$J>0$。基态沿$+z$排列。Holstein--Primakoff精确表示$S^z=S-a^\dagger a$、$S^+=\sqrt{2S-a^\dagger a}\,a$、$S^-=a^\dagger\sqrt{2S-a^\dagger a}$，在低激发密度$\avg{a^\dagger a}\ll2S$时展开根号。保留二次项并Fourier变换得
\begin{equation}
H=E_0+JSz\sum_k(1-\gamma_k)a_k^\dagger a_k,\qquad\gamma_k=\frac1z\sum_\delta\e^{\ii k\cdot\delta}.
\end{equation}
简单立方晶格$\gamma_k\simeq1-a^2k^2/(2d)$，故$\omega_k\propto k^2$。磁化减少$\delta M\propto\int\dd^dk\,n_B(\hbar\omega_k)$，三维给$T^{3/2}$的Bloch律。连续自旋对称下$d\le2$的红外涨落破坏有限温长程序；各向异性或长程作用改变该结论的条件。
\originalsection{反铁磁自旋波与自旋翻转}
双子格反铁磁$H=J\sum_{\avg{ij}}\vct S_i\cdot\vct S_j$，先把B子格自旋坐标旋转，使局域参考态同向，再作HP展开。二次哈密顿量为
\begin{equation}
H_2=JSz\sum_k[a_k^\dagger a_k+b_k^\dagger b_k+\gamma_k(a_kb_{-k}+a_k^\dagger b_{-k}^\dagger)].
\end{equation}
玻色Bogoliubov变换必须满足$u_k^2-v_k^2=1$，于是$\hbar\omega_k=JSz\sqrt{1-\gamma_k^2}$，长波线性。铁磁与反铁磁的Goldstone模色散不同，原因在于低能时间导数结构及未破缺生成元之间的对易，而不是都由“有连续对称”就得相同色散。

加易轴各向异性与沿轴磁场时，Néel态可能转为自旋翻转态：两子格主要反向却横于磁场，并有共同小纵向分量。若定义易轴能密度为$-K\ell_z^2/2$，横向磁化能为$m^2/(2\chi_\perp)-Hm_z$且$\vct m\perp\vct\ell$，极小化$\vct m$给场能$-\chi_\perp H^2(1-\ell_z^2)/2$。比较易轴与横向极小，转变尺度$H_{\rm sf}=\sqrt{K/\chi_\perp}$。系数取决于所定义的$K$，这比不注明哈密顿量就给一个$\sqrt{JK}$更可复核。

中子磁散射测动态自旋结构因子$S^{\alpha\beta}(q,\omega)$。线性自旋波的单磁振子给$\delta(\omega-\omega_q)$峰，多磁振子给连续谱，强度含极化因子与子格干涉。只知道能谱不足以知道某种探针能否看到该模。
\originalsection{巡游磁性与自旋密度波}
在单带Hubbard平均场中，设$n_\uparrow=n/2+m$、$n_\downarrow=n/2-m$。相互作用能$Un_\uparrow n_\downarrow=Un^2/4-Um^2$，自旋分裂带来的动能代价近$ m^2/D_\uparrow(E_F)$。二次系数变负的Stoner条件为$UD_\uparrow(E_F)>1$。采用双自旋总态密度时应写$UD(E_F)/2>1$，避免因约定产生两倍差。

更一般地，密度在波矢$Q$处的自旋不稳定由$1-U\chi_0(Q,0)=0$决定。费米面嵌套若使$\epsilon_{k+Q}\simeq-\epsilon_k$，静态粒子--空穴响应可有大对数，从而优先产生自旋密度波而非均匀铁磁。平均场混合$k\uparrow$与$k+Q\downarrow$，二维矩阵本征能为$\frac12(\epsilon_k+\epsilon_{k+Q})\pm\sqrt{[(\epsilon_k-\epsilon_{k+Q})/2]^2+\Delta_Q^2}$。只有被嵌套部分打开隙时可仍有剩余费米面。平均场忽略强涨落，尤其低维系统不能以Stoner条件当作精确相边界。
